\chapter{2015年福建卷（理）}

\section{单选题}

\begin{problem}
若集合 \(A = \{\i, \i^2, \i^3, \i^4\}\)（\(\i\) 是虚数单位），\(B = \{1, -1\}\)，则 \(A \cap B =\)（\quad）

\choices
    {\(\{-1\}\)}
    {\(\{1\}\)}
    {\(\{1, -1\}\)}
    {\(\varnothing\)}
\end{problem}

\begin{answer}
C
\end{answer}

\begin{solution}
由虚数单位的性质，\(\i^2=-1\)，\(\i^3=-\i\)，\(\i^4=1\)，所以
\(A=\{\i,-1,-\i,1\}=\{1,-1,\i,-\i\}\).因此 \(A\cap B=\{1,-1\}\)，选 C.
\end{solution}

\begin{problem}
下列函数为奇函数的是（\quad）

\choices
    {\( y = \sqrt{x} \)}
    {\( y = |\sin x| \)}
    {\( y = \cos x \)}
    {\( y = \e^{x} - \e^{-x} \)}
\end{problem}

\begin{answer}
D
\end{answer}

\begin{solution}
函数成为奇函数，首先要求定义域关于原点对称，并且满足 \(f(-x)=-f(x)\).函数 \(y=\sqrt{x}\) 的定义域为 \([0,+\infty)\)，不关于原点对称；\(y=|\sin x|\) 和 \(y=\cos x\) 都是偶函数.对 \(f(x)=\e^x-\e^{-x}\)，有
\(f(-x)=\e^{-x}-\e^x=-f(x)\)，故它是奇函数，选 D.
\end{solution}

\begin{problem}
若双曲线 \(E: \frac{x^2}{9} - \frac{y^2}{16} = 1\) 的左、右焦点分别为 \(F_1\)、\(F_2\)，点 \(P\) 在双曲线 \(E\) 上，且 \(|PF_1| = 3\)，则 \(|PF_2|\) 等于（\quad）

\choices
    {11}
    {9}
    {5}
    {3}
\end{problem}

\begin{answer}
B
\end{answer}

\begin{solution}
双曲线的实半轴 \(a=3\)，虚半轴 \(b=4\)，焦半距 \(c=\sqrt{a^2+b^2}=5\)，故焦点为 \(F_1(-5,0)\)、\(F_2(5,0)\).由双曲线定义，点 \(P\) 到两焦点的距离之差的绝对值为 \(2a=6\)，所以
\(\bigl||PF_1|-|PF_2|\bigr|=6\).代入 \(|PF_1|=3\)，得 \(|3-|PF_2||=6\).设 \(d=|PF_2|\geqslant0\)，则 \(|3-d|=6\)，从而 \(d=-3\) 或 \(d=9\)，舍去不符合距离非负性的 \(d=-3\)，故 \(|PF_2|=9\)，选 B.
\end{solution}

\begin{problem}
为了解某社区居民的家庭年收入与年支出的关系，随机调查了该社区 5 户家庭，得到如下统计数据表：

\begin{center}
\begin{tabular}{c|ccccc}
收入 \(x\)（万元） & \(8.2\) & \(8.6\) & \(10.0\) & \(11.3\) & \(11.9\) \\
支出 \(y\)（万元） & \(6.2\) & \(7.5\) & \(8.0\) & \(8.5\) & \(9.8\)
\end{tabular}
\end{center}

根据上表可得回归直线方程 \(\hat{y}=\hat{b}x+\hat{a}\)，其中 \(\hat{b}=0.76\)，\(\hat{a}=\bar{y}-\hat{b}\bar{x}\). 据此估计，该社区一户年收入为 \(15\) 万元家庭的年支出为（\quad）

\choices
    {\(11.4\text{~万元}\)}
    {\(11.8\text{~万元}\)}
    {\(12.0\text{~万元}\)}
    {\(12.2\text{~万元}\)}
\end{problem}

\begin{answer}
B
\end{answer}

\begin{solution}
由表中数据，\(\bar{x}=\dfrac{8.2+8.6+10.0+11.3+11.9}{5}=10\)，\(\bar{y}=\dfrac{6.2+7.5+8.0+8.5+9.8}{5}=8\).因此
\(\hat{a}=\bar{y}-\hat{b}\bar{x}=8-0.76\times10=0.4\).当收入 \(x=15\) 万元时，回归直线给出的支出估计为
\(\hat{y}=0.76\times15+0.4=11.8\)（万元），选 B.
\end{solution}

\begin{problem}
若变量 \(x\)，\(y\) 满足约束条件 \(\begin{cases}
x + 2y \geqslant 0,\\
x - y \leqslant 0,\\
x - 2y + 2 \geqslant 0,
\end{cases}\) 则 \(z = 2x - y\) 的最小值等于（\quad）

\choices
    {\(-\frac{5}{2}\)}
    {\(-2\)}
    {\(-\frac{3}{2}\)}
    {2}
\end{problem}

\begin{answer}
A
\end{answer}

\begin{solution}
三条边界直线分别为 \(y=-\dfrac{x}{2}\)、\(y=x\) 和 \(y=\dfrac{x}{2}+1\).可行域的三个顶点由边界直线两两相交得到：
\((-1,\dfrac12)\)、\((0,0)\)、\((2,2)\).目标函数在多边形可行域的顶点处取得最值，逐点计算得
\(z(-1,\dfrac12)=-2-\dfrac12=-\dfrac52\)，\(z(0,0)=0\)，\(z(2,2)=4-2=2\).所以最小值为 \(-\dfrac52\)，选 A.
\end{solution}

\begin{problem}
阅读如图所示的程序框图，运行相应的程序，则输出的结果为（\quad）

\bitmapfigure[max width=0.52\linewidth]{img/2015/fujian_science/q6_fig1.png}

\choices
    {2}
    {1}
    {0}
    {\(-1\)}
\end{problem}

\begin{answer}
C
\end{answer}

\begin{solution}
循环第一次取 \(i=1\)，加入 \(\cos\dfrac{\pi}{2}=0\)；随后依次加入
\(\cos\pi=-1\)、\(\cos\dfrac{3\pi}{2}=0\)、\(\cos2\pi=1\)、\(\cos\dfrac{5\pi}{2}=0\).此时 \(i\) 增加到 \(6\)，满足 \(i>5\)，循环结束，输出
\(S=0-1+0+1+0=0\)，选 C.
\end{solution}

\begin{problem}
若 \(l\)、\(m\) 是两条不同的直线，\(m\) 垂直于平面 \(\alpha\)，则“\(l \perp m\)”是“\(l \parallel \alpha\)”的（\quad）

\choices
    {充分而不必要条件}
    {必要而不充分条件}
    {充分必要条件}
    {既不充分也不必要条件}
\end{problem}

\begin{answer}
B
\end{answer}

\begin{solution}
因为 \(m\perp\) 平面 \(\alpha\)，所以 \(m\) 垂直于平面 \(\alpha\) 内的任意直线.若 \(l\parallel\alpha\)，则 \(l\) 的方向与平面 \(\alpha\) 内某条直线的方向相同，故 \(l\perp m\).因此 \(l\perp m\) 是 \(l\parallel\alpha\) 的必要条件.

但 \(l\perp m\) 不足以推出 \(l\parallel\alpha\)：例如取 \(l\) 为平面 \(\alpha\) 内且经过 \(m\) 与平面 \(\alpha\) 交点的直线，则 \(l\perp m\)，而 \(l\) 不平行于 \(\alpha\).故选 B.
\end{solution}

\begin{problem}
若 \(a\)、\(b\) 是函数 \(f(x) = x^2 - px + q\)（\(p > 0\)，\(q > 0\)）的两个不同的零点，且 \(a\)、\(b\)、\(-2\) 这三个数可适当排序后成等差数列，也可适当排序后成等比数列，则 \(p + q\) 的值等于（\quad）

\choices
    {6}
    {7}
    {8}
    {9}
\end{problem}

\begin{answer}
D
\end{answer}

\begin{solution}
由于 \(p>0\)、\(q>0\)，且 \(a+b=p>0\)、\(ab=q>0\)，所以两个零点 \(a\)，\(b\) 都为正数.三个数 \(a\)，\(b\)，\(-2\) 成等比数列时，负数 \(-2\) 只能位于中间，因此
\(ab=(-2)^2=4\)，即 \(q=4\).

三个数成等差数列时，\(-2\) 不可能为中项，因为这将导致 \(a+b=-4\).设 \(a\) 为中项，则 \(2a=b-2\)，结合 \(ab=4\) 得 \(a(2a+2)=4\)，即 \((a-1)(a+2)=0\).因 \(a>0\)，得 \(a=1\)，\(b=4\)；交换 \(a\)，\(b\) 只改变两根次序.因此 \(p=a+b=5\)，从而 \(p+q=5+4=9\)，选 D.
\end{solution}

\begin{problem}
已知 \(\overrightarrow{AB} \perp \overrightarrow{AC}\)，\(|\overrightarrow{AB}| = \frac{1}{t}\)，\(|\overrightarrow{AC}| = t\). 若点 \(P\) 是 \(\triangle ABC\) 所在平面内的一点，且 \(\overrightarrow{AP} = \frac{\overrightarrow{AB}}{|\overrightarrow{AB}|} + \frac{4\overrightarrow{AC}}{|\overrightarrow{AC}|}\)，则 \(\overrightarrow{PB} \cdot \overrightarrow{PC}\) 的最大值等于（\quad）

\choices
    {13}
    {15}
    {19}
    {21}
\end{problem}

\begin{answer}
A
\end{answer}

\begin{solution}
由长度条件知 \(t>0\).取 \(A\) 为原点，令 \(\overrightarrow{AB}\) 和 \(\overrightarrow{AC}\) 的方向分别为两条互相垂直的坐标轴正方向.则
\(B=(\dfrac1t,0)\)、\(C=(0,t)\)，而
\(P=(1,4)\).所以
\(\overrightarrow{PB}=\left(\frac1t-1,-4\right)\)，\(\overrightarrow{PC}=(-1,t-4)\).
于是
\[
\overrightarrow{PB}\cdot\overrightarrow{PC}
=-\left(\frac1t-1\right)-4(t-4)
=17-\left(\frac1t+4t\right)
\]
由于 \(t>0\)，由基本不等式
\(\dfrac1t+4t\geqslant2\sqrt{4}=4\)，且当 \(\dfrac1t=4t\)，即 \(t=\dfrac12\) 时取等号.因此最大值为 \(17-4=13\)，选 A.
\end{solution}

\begin{problem}
若定义在 \(\R\) 上的函数 \(f(x)\) 满足 \(f(0) = -1\)，其导函数 \(f'(x)\) 满足 \(f'(x) > k > 1\)，则下列结论中一定错误的是（\quad）

\choices
    {\(f\left(\frac{1}{k}\right) < \frac{1}{k}\)}
    {\(f\left(\frac{1}{k}\right) > \frac{1}{k - 1}\)}
    {\(f\left(\frac{1}{k - 1}\right) < \frac{1}{k - 1}\)}
    {\(f\left(\frac{1}{k - 1}\right) > \frac{k}{k - 1}\)}
\end{problem}

\begin{answer}
C
\end{answer}

\begin{solution}
由 \(f'(x)>k\) 及拉格朗日中值定理，对任意 \(x>0\)，存在 \(\xi\in(0,x)\)，使得
\(f(x)-f(0)=f'(\xi)x>kx\).又 \(f(0)=-1\)，所以
\(f(x)>kx-1\).取 \(x=\dfrac1{k-1}\)，得到
\[
f\left(\frac1{k-1}\right)>\frac{k}{k-1}-1=\frac1{k-1}
\]
因此选项 C 所写的不等式必不成立.

其余选项并非一定错误.例如取 \(f(x)=cx-1\)，其中 \(c>k\).取 \(k<c<k+1\) 时第一个选项成立，取足够大的 \(c\) 时第二、第四个选项均可成立，故只有第三个选项一定错误，选择第三项.
\end{solution}

\section{填空题}

\begin{problem}
\((x + 2)^{5}\) 的展开式中，\(x^{2}\) 的系数等于\fillinblank{}.
\end{problem}

\begin{answer}
80
\end{answer}

\begin{solution}
二项式通项中含 \(x^2\) 的项为从 \(5\) 个因子中选出两个 \(x\)，其余取 \(2\)，故系数为
\(\mathrm C_5^2 2^3=10\times8=80\).
\end{solution}

\begin{problem}
若锐角 \(\triangle ABC\) 的面积为 \(10\sqrt{3}\)，且 \(AB = 5\)，\(AC = 8\)，则 \(BC\) 等于\fillinblank{}.
\end{problem}

\begin{answer}
7
\end{answer}

\begin{solution}
设 \(\angle A\) 为 \(AB\) 与 \(AC\) 的夹角.由三角形面积公式，
\(\dfrac12\cdot5\cdot8\sin A=10\sqrt3\)，所以 \(\sin A=\dfrac{\sqrt3}{2}\).因三角形为锐角三角形，\(A=\dfrac\pi3\)，从而 \(\cos A=\dfrac12\).由余弦定理，
\[
BC^2=AB^2+AC^2-2AB\cdot AC\cos A
=25+64-2\cdot5\cdot8\cdot\frac12=49
\]
因为 \(BC>0\)，所以 \(BC=7\).
\end{solution}

\begin{problem}
如图，点 \(A\) 的坐标为 \((1,0)\)，点 \(C\) 的坐标为 \((2,4)\)，函数 \(f(x) = x^2\). 若在矩形 \(ABCD\) 内随机取一点，则此点取自阴影部分的概率等于\fillinblank{}.

\bitmapfigure[max width=0.36\linewidth]{img/2015/fujian_science/q13_fig1.png}
\end{problem}

\begin{answer}
\(\dfrac{5}{12}\)
\end{answer}

\begin{solution}
矩形的长为 \(2-1=1\)，高为 \(4-0=4\)，面积为 \(4\).阴影部分位于曲线 \(y=x^2\) 上方、矩形内部，故其面积为
\[
4-\int_1^2x^2\,\mathrm{d}x
=4-\left.\frac{x^3}{3}\right|_1^2
=4-\frac73=\frac53
\]
在矩形内均匀取点时，取到阴影部分的概率等于面积比，故
\(P=\dfrac{5/3}{4}=\dfrac5{12}\).
\end{solution}

\begin{problem}
若函数 \(f(x) = \begin{cases}
-x + 6, & x \leqslant 2,\\
3 + \log_a x, & x > 2,
\end{cases}\)（\(a > 0\)，且 \(a \neq 1\)）的值域是 \([4, +\infty)\)，则实数 \(a\) 的取值范围是\fillinblank{}.
\end{problem}

\begin{answer}
\((1,2]\)
\end{answer}

\begin{solution}
当 \(x\leqslant2\) 时，\(f(x)=-x+6\)，其值域为 \([4,+\infty)\).

若 \(0<a<1\)，则 \(\log_a x\) 在 \((2,+\infty)\) 上递减，并且当 \(x\to+\infty\) 时趋于 \(-\infty\)，从而函数值会小于 \(4\)，不符合题意.

若 \(a>1\)，则 \(3+\log_a x\) 在 \((2,+\infty)\) 上递增.为使这一段不出现小于 \(4\) 的函数值，需要
\(3+\log_a2\geqslant4\)，即 \(\log_a2\geqslant1\).由于 \(a>1\)，这等价于 \(a\leqslant2\).结合 \(a\neq1\)，得 \(1<a\leqslant2\).因此 \(a\) 的取值范围为 \((1,2]\).
\end{solution}

\begin{problem}
一个二元码是由 \(0\) 和 \(1\) 组成的数字串 \(x_{1}x_{2}\cdots x_{n}\)（\(n\in \mathbf{N}^{*}\)），其中 \(x_{k}\)（\(k=1\)，\(2\)，\(\dots\)，\(n\)）称为第 \(k\) 位码元. 二元码是通信中常用的码，但在通信过程中有时会发生码元错误（即码元由 \(0\) 变为 \(1\)，或者由 \(1\) 变为 \(0\)）. 已知某种二元码 \(x_{1}x_{2}\cdots x_{7}\) 的码元满足如下校验方程组：\(\begin{cases}
x_{4}\oplus x_{5}\oplus x_{6}\oplus x_{7}=0,\\
x_{2}\oplus x_{3}\oplus x_{6}\oplus x_{7}=0,\\
x_{1}\oplus x_{3}\oplus x_{5}\oplus x_{7}=0,
\end{cases}\)，其中运算 \(\oplus\) 定义为：\(0\oplus 0=0\)，\(0\oplus 1=1\)，\(1\oplus 0=1\)，\(1\oplus 1=0\). 现已知一个这种二元码在通信过程中仅在第 \(k\) 位发生码元错误后变成了 \(1101101\)，那么利用上述校验方程组可判定 \(k\) 等于\fillinblank{}.
\end{problem}

\begin{answer}
5
\end{answer}

\begin{solution}
将通信后得到的码字 \(1101101\) 代入三个校验式，得到校验结果
\(s_1=1\oplus1\oplus0\oplus1=1\)，\(s_2=1\oplus0\oplus0\oplus1=0\)，\(s_3=1\oplus0\oplus1\oplus1=1\).
若第 \(k\) 位发生一次错误，则包含该位的校验式结果由 \(0\) 变为 \(1\)，不包含该位的校验式结果仍为 \(0\).各位对应的校验结果依次为
\((0,0,1)\)，\((0,1,0)\)，\((0,1,1)\)，\((1,0,0)\)，\((1,0,1)\)，\((1,1,0)\)，\((1,1,1)\).
观测到 \((s_1,s_2,s_3)=(1,0,1)\)，对应第 \(5\) 位，所以 \(k=5\).
\end{solution}

\section{解答题}

\begin{problem}
某银行规定，一张银行卡若在一天内出现 \(3\) 次密码尝试错误，该银行卡将被锁定. 小王到该银行取钱时，发现自己忘记了银行卡的密码，但可以确认该银行卡的正确密码是他常用的 \(6\) 个密码之一，小王决定从中不重复地随机选择 \(1\) 个进行尝试. 若密码正确，则结束尝试；否则继续尝试，直至该银行卡被锁定.
\begin{enumerate}
\item 求当天小王的该银行卡被锁定的概率；
\item 设当天小王用该银行卡尝试密码的次数为 \(X\)，求 \(X\) 的分布列和数学期望.
\end{enumerate}
\end{problem}

\begin{solution}
设正确密码在随机排列中的位置为 \(R\).由于六个密码等可能地被排列，正确密码出现在每个位置的概率都是 \(\dfrac16\).

\begin{enumerate}
\item 银行卡在前三次尝试中都错误时被锁定，这等价于 \(R\geqslant4\).因此
\[
P(\text{银行卡被锁定})=\frac{3}{6}=\frac12
\]

\item 随机变量 \(X\) 的取值为 \(1\)，\(2\)，\(3\).当正确密码在第1次或第2次出现时，分别有
\(P(X=1)=\frac16\)，\(P(X=2)=\frac16\).
当正确密码在第3次、第4次、第5次或第6次出现时，小王都会进行3次尝试，所以
\[
P(X=3)=\frac46=\frac23
\]
因此分布列为：\(P(X=1)=\dfrac16\)、\(P(X=2)=\dfrac16\)、\(P(X=3)=\dfrac23\).
数学期望为
\[
E(X)=1\cdot\frac16+2\cdot\frac16+3\cdot\frac23=\frac52
\]
\end{enumerate}
\end{solution}

\begin{problem}
如图，在几何体 \(ABCDE\) 中，四边形 \(ABCD\) 是矩形，\(AB \perp\) 平面 \(BEC\)，\(BE \perp EC\)，\(AB = BE = EC = 2\)，\(G\)、\(F\) 分别是线段 \(BE\)、\(DC\) 的中点.
\begin{enumerate}
\item 求证：\(GF \parallel\) 平面 \(ADE\)；
\item 求平面 \(AEF\) 与平面 \(BEC\) 所成锐二面角的余弦值.
\end{enumerate}
\bitmapfigure[max width=0.38\linewidth]{img/2015/fujian_science/q17_fig1.png}
\end{problem}

\begin{solution}
建立空间直角坐标系，取 \(B\) 为原点，\(BE\) 所在方向、\(EC\) 所在方向、\(BA\) 的方向分别为三个坐标轴正方向.由 \(BE\perp EC\)、\(AB\perp\) 平面 \(BEC\) 及已知长度均为 \(2\)，可取
\(B=(0,0,0)\)，\(E=(2,0,0)\)，\(C=(2,2,0)\)，\(A=(0,0,2)\).
因为 \(ABCD\) 是矩形，所以 \(D=(2,2,2)\).又 \(G\)、\(F\) 分别为 \(BE\)、\(DC\) 的中点，故
\(G=(1,0,0)\)，\(F=(2,2,1)\).

\begin{enumerate}
\item 平面 \(ADE\) 内有
\(\overrightarrow{AD}=(2,2,0)\)、\(\overrightarrow{AE}=(2,0,-2)\)，且
\(\overrightarrow{AD}\times\overrightarrow{AE}=(-4,4,-4)\)，故可取平面 \(ADE\) 的一个法向量为 \(\overrightarrow{n}=(1,-1,1)\).而
\(\overrightarrow{GF}=(1,2,1)\)，满足
\(\overrightarrow{GF}\cdot\overrightarrow{n}=1-2+1=0\)，所以 \(GF\) 的方向与平面 \(ADE\) 平行.
此外，点 \(G\) 不在平面 \(ADE\) 上，因为平面 \(ADE\) 的方程为 \(x-y+z-2=0\)，而 \(G\) 的坐标满足 \(x-y+z-2=-1\).因此 \(GF\parallel\) 平面 \(ADE\).

\item 平面 \(BEC\) 为 \(z=0\)，其一个法向量为 \(\overrightarrow{n_2}=(0,0,1)\).平面 \(AEF\) 内有
\(\overrightarrow{AE}=(2,0,-2)\)、\(\overrightarrow{AF}=(2,2,-1)\)，故其一个法向量为
\[
\overrightarrow{n_1}=\overrightarrow{AE}\times\overrightarrow{AF}=(4,-2,4)
\]
平面所成锐二面角的余弦等于两法向量夹角余弦的绝对值，因此
\[
\cos\theta=\frac{|\overrightarrow{n_1}\cdot\overrightarrow{n_2}|}{|\overrightarrow{n_1}|\,|\overrightarrow{n_2}|}
=\frac{4}{\sqrt{4^2+2^2+4^2}}=\frac{4}{6}=\frac23
\]
\end{enumerate}
\end{solution}

\begin{problem}
已知椭圆 \(E: \frac{x^2}{a^2} + \frac{y^2}{b^2} = 1\)（\(a > b > 0\)）过点 \((0, \sqrt{2})\)，且离心率 \(e = \frac{\sqrt{2}}{2}\).
\begin{enumerate}
\item 求椭圆 \(E\) 的方程；
\item 设直线 \(l: x = my - 1\)（\(m \in \R\)）交椭圆 \(E\) 于 \(A\)、\(B\) 两点，判断点 \(G\left(-\frac{9}{4}, 0\right)\) 与以线段 \(AB\) 为直径的圆的位置关系，并说明理由.
\end{enumerate}
\bitmapfigure[max width=0.52\linewidth]{img/2015/fujian_science/q18_fig1.png}
\end{problem}

\begin{solution}
\begin{enumerate}
\item 点 \((0,\sqrt2)\) 在椭圆上，所以 \(\dfrac{2}{b^2}=1\)，即 \(b^2=2\).又离心率满足
\(e^2=\dfrac{c^2}{a^2}=\dfrac12\)，而 \(b^2=a^2-c^2\)，故 \(b^2=\dfrac{a^2}{2}\)，从而 \(a^2=4\).因此椭圆方程为
\[
\frac{x^2}{4}+\frac{y^2}{2}=1
\]

\item 设直线 \(l\) 与椭圆交点 \(A\)、\(B\) 的纵坐标分别为 \(y_1\)、\(y_2\).由 \(x=my-1\) 代入椭圆方程，得
\[
(m^2+2)y^2-2my-3=0
\]
所以
\(y_1+y_2=\frac{2m}{m^2+2}\)，\(y_1y_2=-\frac{3}{m^2+2}\).
点 \(G=(-\dfrac94,0)\)，且交点坐标为 \((my_i-1,y_i)\).于是
\[
\begin{gathered}
\overrightarrow{GA}\cdot\overrightarrow{GB}
 =(my_1+\tfrac54)(my_2+\tfrac54)+y_1y_2\\
 =(m^2+1)y_1y_2+\frac{5m}{4}(y_1+y_2)+\frac{25}{16}
 =\frac{17m^2+2}{16(m^2+2)}>0.
\end{gathered}
\]
以 \(AB\) 为直径的圆上点满足 \(\overrightarrow{XA}\cdot\overrightarrow{XB}=0\)，圆内点对应数量积小于0，圆外点对应数量积大于0.因此 \(G\) 在该圆外.
\end{enumerate}
\end{solution}

\begin{problem}
已知函数 \(f(x) = \ln (1 + x)\)，\(g(x) = kx\)（\(k \in \R\)）.
\begin{enumerate}
\item 证明：当 \(x > 0\) 时，\(f(x) < x\)；
\item 证明：当 \(k < 1\) 时，存在 \(x_0 > 0\)，使得对任意的 \(x \in (0, x_0)\)，恒有 \(f(x) > g(x)\)；
\item 确定 \(k\) 的所有可能取值，使得存在 \(t > 0\)，对任意的 \(x \in (0, t)\)，恒有 \(|f(x) - g(x)| < x^2\).
\end{enumerate}
\end{problem}

\begin{solution}
\begin{enumerate}
\item 令 \(h(x)=x-\ln(1+x)\)，则 \(h(0)=0\)，且
\(h'(x)=1-\dfrac1{1+x}=\dfrac{x}{1+x}>0\)（\(x>0\)）.所以 \(h(x)>h(0)=0\)，即 \(\ln(1+x)<x\).

\item 令 \(q(x)=\ln(1+x)-\dfrac{x}{1+x}\).有 \(q(0)=0\)，且
\(q'(x)=\dfrac{x}{(1+x)^2}>0\)（\(x>0\)），所以 \(\ln(1+x)>\dfrac{x}{1+x}\).

若 \(k\leqslant0\)，则对所有 \(x>0\)，有 \(\dfrac{x}{1+x}>kx\)，任取一个正数作为 \(x_0\) 即可.若 \(0<k<1\)，取 \(x_0=\dfrac{1-k}{k}>0\).当 \(0<x<x_0\) 时，\(\dfrac1{1+x}>k\)，故
\[
f(x)=\ln(1+x)>\frac{x}{1+x}>kx=g(x)
\]
因此对任意 \(k<1\)，都存在符合条件的 \(x_0>0\).

\item 先证必要性.由 \(\ln(1+x)<x\)，当 \(k>1\) 时
\[
kx-\ln(1+x)>(k-1)x
\]
若题设中的 \(t>0\) 存在，取 \(0<x<\min\{t,k-1\}\)，则 \((k-1)x>x^2\)，从而 \(|f(x)-g(x)|>x^2\)，矛盾.

当 \(k<1\) 时，令
\(r(x)=\ln(1+x)-x+\dfrac{x^2}{2}\).有 \(r(0)=0\)，且
\(r'(x)=\dfrac{x^2}{1+x}>0\)，所以 \(\ln(1+x)>x-\dfrac{x^2}{2}\).记 \(d=1-k>0\)，则
\[
f(x)-g(x)>dx-\frac{x^2}{2}
\]
若题设中的 \(t>0\) 存在，取 \(0<x<\min\{t,\dfrac{2d}{3}\}\)，则右端大于 \(x^2\)，即 \(|f(x)-g(x)|>x^2\)，矛盾.因此必须有 \(k=1\).

当 \(k=1\) 时，由
\(x-\ln(1+x)=\displaystyle\int_0^x\frac{s}{1+s}\,\mathrm{d}s\)，且 \(0<\dfrac{s}{1+s}<s\)，可得
\[
0<x-\ln(1+x)<\int_0^x s\,\mathrm{d}s=\frac{x^2}{2}<x^2
\]
所以任取 \(t>0\) 都满足 \(|f(x)-g(x)|<x^2\).综上，\(k=1\).
\end{enumerate}
\end{solution}

\begin{problem}
在平面直角坐标系 \(xOy\) 中，圆 \(C\) 的参数方程为 \(\begin{cases}
x = 1 + 3\cos t,\\
y = -2 + 3\sin t.
\end{cases}\)（\(t\) 为参数）. 在极坐标系（与平面直角坐标系 \(xOy\) 取相同的长度单位，且以原点 \(O\) 为极点，以 \(x\) 轴非负半轴为极轴）中，直线 \(l\) 的方程为 \(\sqrt{2}\rho\sin \left(\theta - \frac{\pi}{4}\right) = m\)（\(m \in \R\)）.
\begin{enumerate}
\item 求圆 \(C\) 的普通方程及直线 \(l\) 的直角坐标方程；
\item 设圆心 \(C'\) 到直线 \(l\) 的距离等于 2，求 \(m\) 的值.
\end{enumerate}
\end{problem}

\begin{solution}
\begin{enumerate}
\item 由参数方程消去参数 \(t\)，得
\[
(x-1)^2+(y+2)^2=9
\]
又由极坐标与直角坐标的关系 \(x=\rho\cos\theta\)、\(y=\rho\sin\theta\)，直线方程化为
\[
\sqrt2\rho\sin\left(\theta-\frac\pi4\right)
=\rho(\sin\theta-\cos\theta)=y-x=m
\]
所以直线 \(l\) 的直角坐标方程为 \(x-y+m=0\).

\item 圆心为 \(C_0=(1,-2)\).由点到直线的距离公式，
\[
\frac{|1-(-2)+m|}{\sqrt{1^2+(-1)^2}}=2
\]
即 \(|m+3|=2\sqrt2\).因此
\[
m=-3+2\sqrt2\quad\text{或}\quad m=-3-2\sqrt2
\]
\end{enumerate}
\end{solution}

\section{选考题：解答题}

\begin{problem}
已知函数 \(f(x)\) 的图象是由函数 \(g(x) = \cos x\) 的图象经如下变换得到：先将 \(g(x)\) 图象上所有点的纵坐标伸长到原来的 \(2\) 倍（横坐标不变），再将所得到的图象向右平移 \(\frac{\pi}{2}\) 个单位长度.
\begin{enumerate}
\item 求函数 \(f(x)\) 的解析式，并求其图象的对称轴方程；
\item 已知关于 \(x\) 的方程 \(f(x)+g(x)=m\) 在 \([0,2\pi)\) 内有两个不同的解 \(\alpha\)、\(\beta\).
\begin{enumerate}
\item 求实数 \(m\) 的取值范围；
\item 证明：\(\cos(\alpha-\beta)=\frac{2m^{2}}{5}-1\).
\end{enumerate}
\end{enumerate}
\end{problem}

\begin{solution}
\begin{enumerate}
\item 纵坐标伸长为原来的2倍得到 \(y=2\cos x\)，再向右平移 \(\dfrac\pi2\)，所以
\[
f(x)=2\cos\left(x-\frac\pi2\right)=2\sin x
\]
函数 \(y=2\sin x\) 在取得最大值或最小值处关于竖直直线对称，因此对称轴为
\(x=\dfrac\pi2+k\pi\)，其中 \(k\in\Z\).

\item
\begin{enumerate}
\item 方程化为
\[
2\sin x+\cos x=m
\]
令 \(\cos\varphi=\dfrac2{\sqrt5}\)、\(\sin\varphi=\dfrac1{\sqrt5}\)，则
\[
2\sin x+\cos x=\sqrt5\sin(x+\varphi)
\]
在长度为一个周期的区间 \([0,2\pi)\) 内，方程有两个不同解，当且仅当右端常数严格介于最大值和最小值之间，即
\[
-\sqrt5<m<\sqrt5
\]

\item 设 \(\theta_1=\alpha+\varphi\)、\(\theta_2=\beta+\varphi\)，则
\(\sin\theta_1=\sin\theta_2=\dfrac{m}{\sqrt5}\)，且两解对应的角满足 \(\theta_1+\theta_2\equiv\pi\pmod{2\pi}\).因此
\[
\begin{gathered}
\cos(\alpha-\beta)
 =\cos(\theta_1-\theta_2)=\cos(2\theta_1-\pi)\\
 =-\cos2\theta_1=2\sin^2\theta_1-1=\frac{2m^2}{5}-1.
\end{gathered}
\]
\end{enumerate}
\end{enumerate}
\end{solution}

\begin{problem}
已知矩阵 \(A=\begin{pmatrix}2&1\\4&3\end{pmatrix}\)，\(B=\begin{pmatrix}1&1\\0&-1\end{pmatrix}\).
\begin{enumerate}
\item 求 \(A\) 的逆矩阵 \( A^{-1} \)；
\item 求矩阵 \(C\)，使得 \(AC = B\).
\end{enumerate}
\end{problem}

\begin{solution}
\begin{enumerate}
\item \(\det A=2\cdot3-1\cdot4=2\neq0\)，故 \(A\) 可逆，且
\[
A^{-1}=\frac1{2}\begin{pmatrix}3&-1\\-4&2\end{pmatrix}
=\begin{pmatrix}\frac32&-\frac12\\-2&1\end{pmatrix}
\]

\item 由 \(AC=B\)，左乘 \(A^{-1}\) 得 \(C=A^{-1}B\).因此
\[
C=\begin{pmatrix}\frac32&-\frac12\\-2&1\end{pmatrix}
\begin{pmatrix}1&1\\0&-1\end{pmatrix}
=\begin{pmatrix}\frac32&2\\-2&-3\end{pmatrix}
\]
直接相乘可得 \(AC=B\)，故该矩阵满足条件.
\end{enumerate}
\end{solution}

\begin{problem}
已知 \(a > 0\)，\(b > 0\)，\(c > 0\)，函数 \(f(x) = |x + a| + |x - b| + c\) 的最小值为 \(4\).
\begin{enumerate}
\item 求 \(a + b + c\) 的值；
\item 求 \(\frac{1}{4}a^{2}+\frac{1}{9}b^{2}+c^{2}\) 的最小值.
\end{enumerate}
\end{problem}

\begin{solution}
\begin{enumerate}
\item 当 \(-a\leqslant x\leqslant b\) 时，
\(|x+a|+|x-b|=(x+a)+(b-x)=a+b\).在区间外，该和分别随 \(x\) 远离区间而增大，所以函数最小值为 \(a+b+c\).由题设最小值为4，得
\[
a+b+c=4
\]

\item 由柯西不等式，
\[
(a+b+c)^2=\left(2\cdot\frac a2+3\cdot\frac b3+1\cdot c\right)^2
\leqslant(2^2+3^2+1^2)\left(\frac{a^2}{4}+\frac{b^2}{9}+c^2\right)
\]
代入 \(a+b+c=4\)，得
\[
\frac{a^2}{4}+\frac{b^2}{9}+c^2\geqslant\frac{16}{14}=\frac87
\]
当且仅当 \(\dfrac a2:\dfrac b3:c=2:3:1\) 时取等号，即 \(a:b:c=4:9:1\).结合 \(a+b+c=4\)，可取
\(a=\dfrac87\)、\(b=\dfrac{18}{7}\)、\(c=\dfrac27\)，均为正数，故等号可以取得，最小值为 \(\dfrac87\).
\end{enumerate}
\end{solution}
